\section{Testing \& Continuous Delivery}

% --- SLIDE 1: WHY ML TESTING IS DIFFERENT ---
\begin{frame}{Why You Cannot Just Write Unit Tests}

    \centering
    \begin{columns}[T]
        \begin{column}{0.48\textwidth}
            \begin{tcolorbox}[colback=blue!5, colframe=blue!50, title=\small Ordinary software]
                \small
                \begin{itemize}
                    \item The correct output is \textbf{specified in advance}.
                    \item \texttt{assert add(2,2) == 4}
                    \item A test failing means the code is wrong.
                    \item Behaviour changes only when code changes.
                \end{itemize}
            \end{tcolorbox}
        \end{column}
        \begin{column}{0.48\textwidth}
            \begin{tcolorbox}[colback=red!5, colframe=red!50, title=\small Machine learning]
                \small
                \begin{itemize}
                    \item The correct output is \textbf{learned}, and only known statistically.
                    \item \texttt{assert model.predict(x) == ?}
                    \item A metric dropping might mean the code is wrong, the data changed, or the seed
                          differed.
                    \item Behaviour changes when \emph{data} changes, with no commit at all.
                \end{itemize}
            \end{tcolorbox}
        \end{column}
    \end{columns}

    \vspace{0.9em}
    \begin{block}{So you test three separate things}
        \textbf{(1)} the \emph{code} around the model, with ordinary unit tests;
        \textbf{(2)} the \emph{data} going in, with schema and quality checks;
        \textbf{(3)} the \emph{model's behaviour}, with properties and thresholds rather than exact values.
    \end{block}

\end{frame}
